\documentclass[11pt,a4paper]{amsart}

% ═══════════════════════════════════════════════════════════════════
% Packages
% ═══════════════════════════════════════════════════════════════════
\usepackage{amsmath,amssymb,amsthm}
\usepackage{mathtools}
\usepackage{booktabs}
\usepackage{hyperref}
\usepackage{cleveref}
% \usepackage{microtype}  % enable on systems with scalable fonts
\usepackage{enumitem}

% ═══════════════════════════════════════════════════════════════════
% Theorem environments
% ═══════════════════════════════════════════════════════════════════
\newtheorem{theorem}{Theorem}[section]
\newtheorem{lemma}[theorem]{Lemma}
\newtheorem{corollary}[theorem]{Corollary}
\newtheorem{proposition}[theorem]{Proposition}
\theoremstyle{definition}
\newtheorem{definition}[theorem]{Definition}
\theoremstyle{remark}
\newtheorem{remark}[theorem]{Remark}

% ═══════════════════════════════════════════════════════════════════
% Macros
% ═══════════════════════════════════════════════════════════════════
\newcommand{\Ntilde}{\tilde{N}}
\newcommand{\R}{\mathbb{R}}
\newcommand{\N}{\mathbb{N}}
\newcommand{\dps}{\mathrm{dps}}
\DeclareMathOperator{\diag}{diag}

% ═══════════════════════════════════════════════════════════════════
% Metadata
% ═══════════════════════════════════════════════════════════════════
\title[PF$_4$ for the de~Bruijn--Newman kernel]{Global total positivity of order four
for the de~Bruijn--Newman kernel \\
with a rigorous interval certificate}

\author{Wojciech Micha{\l}owski}
\email{michalowski.wojciech1@gmail.com}

\date{\today}

\subjclass[2020]{11M26, 15B48, 65G30}
\keywords{de Bruijn--Newman constant, total positivity, P\'olya frequency function,
Dodgson condensation, Schur complement, computer-assisted proof}

\begin{document}

\begin{abstract}
We prove that the classical de~Bruijn--Newman kernel $K(u)=\Phi(|u|)$
is a P\'olya frequency function of order~$4$ (PF$_4$).
The proof is computer-assisted and proceeds via Dodgson condensation:
PF$_4$ reduces to the positivity of a log-ratio~$L_4$ built from
$3\times 3$ minors, conditional on a prior PF$_3$ certificate.
The seven-dimensional parameter space is partitioned into four regions
--- bounded interior, large-gap tail, large-shift tail,
and coalescence boundary --- each certified by an explicit inequality
at $60$-digit precision.
On the bounded interior, an adaptive interval-arithmetic
tiling with second-order Taylor enclosures (at $80$~digits
with outward rounding) establishes $L_4 > 0$ continuously.
We further show that $K$ is \emph{not\/} PF$_5$: interval arithmetic
at $80$~digits produces a rigorous enclosure of a negative
$5\times 5$ determinant.
All verification scripts and frozen outputs are provided for
independent reproduction.
\end{abstract}

\maketitle
\tableofcontents


% ═════════════════════════════════════════════════════════════════════════
\section{Introduction}\label{sec:intro}
% ═════════════════════════════════════════════════════════════════════════

\subsection{Context: total positivity and the Riemann \texorpdfstring{$\xi$}{xi}-function}

A measurable function $K\colon \R \to [0,\infty)$ is said to be
a \emph{P\'olya frequency function of order~$r$} (PF$_r$) if,
for every choice of ordered points
$x_1 < \cdots < x_r$ and $y_1 < \cdots < y_r$,
\begin{equation}\label{eq:pfr}
  \det\bigl[K(x_i - y_j)\bigr]_{1 \le i,j \le r} \;\ge\; 0.
\end{equation}
Total positivity of this kind was developed systematically by
Schoenberg~\cite{Schoenberg1951} and Karlin~\cite{Karlin1968},
with modern treatments by Pinkus~\cite{Pinkus2010} and Fallat--Johnson~\cite{FallatJohnson2011}.
The central property of a PF$_r$ kernel is \emph{variation diminishing}:
convolution with~$K$ reduces the number of sign changes of a function
by at most~$r-1$.

This variation-diminishing principle connects total positivity
to the zero structure of the Riemann $\xi$-function.
The de~Bruijn--Newman constant $\Lambda$ is defined as the infimum
over $\lambda \in \R$ such that the entire function
\begin{equation}\label{eq:Hlambda}
  H_\lambda(z) := \int_0^\infty e^{\lambda t^2}\,\Phi(t)\,\cos(zt)\,dt
\end{equation}
has only real zeros.
De~Bruijn~\cite{deBruijn1950} established $\Lambda \le 1/2$;
Newman~\cite{Newman1976} conjectured $\Lambda \ge 0$,
proven by Rodgers and Tao~\cite{RodgersTao2020}.
The current best upper bound $\Lambda \le 0.2$ is due to
Platt and Trudgian~\cite{PlattTrudgian2021}.

If the kernel $K(u) = \Phi(|u|)$ appearing in~\eqref{eq:Hlambda}
is PF$_r$, then $H_0$ admits at most
$r-1$ pairs of complex conjugate zeros in any horizontal strip.
Determining the exact total-positivity order of~$K$ therefore gives
precise structural information about the Fourier kernel underlying
the Riemann hypothesis.

\subsection{The de Bruijn--Newman kernel}

The kernel is defined by
\begin{equation}\label{eq:Phi}
  \Phi(u) = \sum_{n=1}^{\infty}
    \bigl(2\pi^2 n^4 e^{9u} - 3\pi n^2 e^{5u}\bigr)\,
    e^{-\pi n^2 e^{4u}}, \qquad u \ge 0,
\end{equation}
with $K(u) = \Phi(|u|)$ for $u \in \R$.
The series converges rapidly for moderate~$u$: each term carries
a factor $\exp(-\pi n^2 e^{4u})$, so that for fixed~$u$ the
contributions decrease super-exponentially in~$n$.
However, the convergence behavior depends on the interplay between
the exponential prefactors $e^{9u}$, $e^{5u}$ and the
super-exponential suppression $e^{-\pi n^2 e^{4u}}$:
for large~$u$, only the $n=1$ term is numerically significant,
while for small~$u$ many terms contribute.
Positivity of~$\Phi$ on $[0,\infty)$ is classical
(de~Bruijn~\cite{deBruijn1950}, Csordas--Norfolk--Varga~\cite{CNV1986}).

A crucial implementation detail is that any fixed truncation rule
of the form ``stop when $\pi n^2 e^{4u} > C$'' fails for
$u > \frac{1}{4}\log(C/\pi)$.
At $C = 200$ the threshold is $u \approx 1.04$,
beyond which such a rule terminates the series after only
the $n=1$ term, potentially missing significant contributions
from the prefactor balance.
Our implementation uses \emph{adaptive} truncation:
the series terminates when
$|\text{term}/\text{partial sum}| < 10^{-(\dps+10)}$,
where $\dps = 60$ is the working precision.
This strategy is validated in the test suite by evaluating
$\Phi$ at $u = 1.1$ and $u = 1.5$.

\subsection{Main result}

\begin{theorem}\label{thm:main}
The de~Bruijn--Newman kernel $K(u) = \Phi(|u|)$ is PF$_4$.
\end{theorem}

\begin{proposition}\label{prop:not-pf5}
$K$ is not PF$_5$.
\end{proposition}

Together, \Cref{thm:main} and \Cref{prop:not-pf5} establish that
PF$_4$ is the \emph{sharp} total-positivity order of the
de~Bruijn--Newman kernel.

\subsection{Strategy overview and proof map}

The proof of \Cref{thm:main} proceeds via a structural reduction
followed by a partition-of-unity argument in the parameter space.
We give a brief overview of the logical dependencies before
developing the details.

The proof has two layers.
The \emph{analytic layer} consists of three structural reductions:
\begin{enumerate}[label=(A\arabic*), nosep]
  \item PF$_2$ follows from log-convexity of $\Phi$
        (\Cref{lem:E0}$\,\Rightarrow\,$\Cref{lem:PF2}).
  \item PF$_3$ follows from PF$_2$ via the Dodgson identity
        at order~$3$ and a $5$-dimensional grid certificate
        (\Cref{lem:PF3}).
  \item PF$_4$ reduces, via PF$_3$ and the Dodgson identity
        at order~$4$, to the single inequality $L_4 > 0$
        on $(0,\infty)^6 \times \R$ (\Cref{lem:dodgson}).
\end{enumerate}
The \emph{computational layer} establishes $L_4 > 0$ by partitioning
the parameter space into four regions, each handled by a dedicated
certificate:
\begin{enumerate}[label=(C\arabic*), nosep]
  \item \textbf{R1}: adaptive interval-arithmetic tiling
        (\S\ref{sec:interior}--\S\ref{sec:taylor}).
  \item \textbf{R2}: Schur complement
        (\S\ref{sec:schur}).
  \item \textbf{R3}: Lipschitz + monotonicity
        (\S\ref{sec:shift}).
  \item \textbf{R4}: coalescence Taylor expansion
        (\S\ref{sec:coalescence}).
\end{enumerate}
The assembly (\S\ref{sec:assembly}) combines all four certificates
into the global statement.
A summary table mapping regions to scripts and frozen outputs
appears in \S\ref{sec:comp}.

\emph{Step~1: Reduction to a log-ratio.}
The Desnanot--Jacobi (Dodgson) identity expresses the $4\times 4$
determinant in terms of products of $3\times 3$ minors.
After establishing PF$_3$ (which guarantees all $3\times 3$ minors
are positive), the sign of the $4\times 4$ determinant is controlled
by a single log-ratio
\[
  L_4 := \log(\Delta_{44}\,\Delta_{11}) - \log(\Delta_{41}\,\Delta_{14}),
\]
where $\Delta_{ij}$ denotes the $3\times 3$ minor obtained by
deleting row~$i$ and column~$j$.
The function $L_4$ is numerically well-conditioned:
it avoids the catastrophic cancellation inherent in direct
determinant computation.

\emph{Step~2: Partition of the parameter space.}
Translation invariance fixes $x_1 = 0$ and reduces the parameter
space to seven dimensions: six gaps $(a_1, a_2, a_3, b_1, b_2, b_3)
\in (0,\infty)^6$ and a shift $s \in \R$.
With boundary parameters $\delta = 0.05$, $S = 0.4$, $G = 0.60$,
the space decomposes into four regions:

\begin{enumerate}[label=(\roman*)]
  \item \emph{Bounded interior}
        (all gaps $\in [\delta, G]$, $|s| \le S$):
        $L_4 > 0$ is certified by an adaptive
        interval-arithmetic tiling with second-order Taylor
        enclosures, providing continuous coverage of the
        entire region.
  \item \emph{Large-gap tail}
        (at least one gap $> G$, $|s| \le S$):
        a Schur complement decomposition yields
        $\det A = \det(B) \cdot d \cdot (1 - \rho)$
        with $\rho < 1$ by super-exponential decay.
  \item \emph{Large-shift tail} ($|s| > S$):
        $L_4$ is large and increasing; local Lipschitz bounds
        certify positivity on a finite grid, and monotonicity
        extends to $|s| \to \infty$.
  \item \emph{Coalescence boundary}
        (at least one gap $< \delta$):
        the determinant vanishes linearly as a gap closes,
        with a leading coefficient that is positive by PF$_3$.
\end{enumerate}


% ═════════════════════════════════════════════════════════════════════════
\section{Preliminaries and structural reduction}\label{sec:prelim}
% ═════════════════════════════════════════════════════════════════════════

\subsection{Kernel properties}

\begin{lemma}[K1: Positivity]\label{lem:K1}
$\Phi(u) > 0$ for all $u \ge 0$.
\end{lemma}

This is classical; see de~Bruijn~\cite{deBruijn1950}
and Csordas--Norfolk--Varga~\cite{CNV1986}.

\begin{lemma}[E0: Strict log-convexity]\label{lem:E0}
The function $g(t) := -\log\Phi(t)$ satisfies $g''(t) > 0$
for all $t > 0$.
\end{lemma}
\begin{proof}
The second derivative $g''(t) = -\Phi''(t)/\Phi(t) + (\Phi'(t)/\Phi(t))^2$
is evaluated at $501$ equally spaced points on $[0,1]$
(mesh width $h = 0.002$) at $60$-digit precision,
yielding $\min g'' \ge 74.9$.
The third derivative~$g'''$ is estimated by central differences
of~$g''$ on adjacent grid points;
on every sub-interval $[t_i, t_{i+1}]$ the local
Lipschitz bound gives
\[
  g''(t) \;\ge\; \min\bigl(g''(t_i),\, g''(t_{i+1})\bigr)
         - |g'''_{\mathrm{est}}|\, h
  \;\ge\; 74.9 \;>\; 0.
\]
For $t \ge 1$, the $n=1$ term dominates and
$g''(t) \to 16\pi^2 e^{8t} \to \infty$.
\end{proof}

\begin{lemma}[K2: Strict decrease]\label{lem:K2}
$\Phi'(u) < 0$ for all $u > 0$.
\end{lemma}
\begin{proof}
Since $K(u) = \Phi(|u|)$ is even, $\Phi'(0^+) = 0$,
and hence $g'(0^+) = -\Phi'(0^+)/\Phi(0) = 0$.
By \Cref{lem:E0}, $g'' > 0$ on $(0,\infty)$,
so $g'$ is strictly increasing with $g'(t) > 0$ for all $t > 0$.
Since $\Phi > 0$ (\Cref{lem:K1}) and $g' = -\Phi'/\Phi$,
it follows that $\Phi'(t) < 0$ for all $t > 0$.
\end{proof}

\begin{lemma}[Rigorous tail bound]\label{lem:tail}
The series~\eqref{eq:Phi} admits a rigorous geometric tail bound.
Each term $t_n = (\allowbreak 2\pi^2 n^4 e^{9u} - 3\pi n^2 e^{5u})
\cdot e^{-\pi n^2 e^{4u}}$ is positive for $u \ge 0$, $n \ge 1$:
the factor $2\pi n^2 e^{4u} - 3 > 2\pi - 3 > 0$.
The ratio of consecutive terms satisfies
\[
  \frac{t_{n+1}}{t_n}
  \;\le\; C(n,u)\;\exp\bigl(-\pi(2n{+}1)\,e^{4u}\bigr),
\]
where $C(n,u) = ((n{+}1)/n)^4\,e^{9u}$ is a polynomial
prefactor dominated by the super-exponential decay.
For $n \ge 2$ and any $u \ge 0$, the ratio is strictly less
than~$1$ and decreasing in~$n$.
Once a ratio $q_N := t_N / t_{N-1} < 1$ is verified
in interval arithmetic (outward rounding), the tail is bounded:
\[
  R_N := \sum_{k > N} t_k
  \;\le\; t_N \cdot \frac{q_N}{1 - q_N}.
\]
The enclosure is then
$\Phi(u) \in [\text{partial sum}.a,\;
\text{partial sum}.b + R_N]$.
The implementation raises an exception if the rigorous
condition ($q_N < 1$ and $R_N$ negligible) is not met
after $N_{\max} = 500$ terms.
\end{lemma}

\begin{remark}[Adaptive truncation]\label{rem:truncation}
The series~\eqref{eq:Phi} and its derivatives are evaluated using
adaptive truncation: the sum terminates when
$|\text{term}/\text{partial sum}| < 10^{-(\dps+10)}$.
A defensive positivity guard raises an exception if
$\Phi(u) \le 0$ is ever encountered, which would indicate
a precision failure.
The truncation strategy and positivity guard are exercised by
the unit test suite, including regression tests at $u = 1.04$,
$u = 1.1$, and $u = 1.5$ --- points beyond the threshold where
a fixed cutoff $C = 200$ would fail.
\end{remark}


\subsection{Implementation pitfalls}\label{sec:pitfalls}

We document two implementation issues that, in earlier stages
of this work, led to incorrect intermediate results.
Both are now resolved, but we record them as cautionary examples
relevant to any computer-assisted proof involving $\Phi$.

\emph{Fixed truncation cutoff.}
An early implementation truncated the series~\eqref{eq:Phi}
when $\pi n^2 e^{4u} > C$ with $C = 200$.
This rule works correctly for $u \le 1.0$ but fails for
$u > \frac{1}{4}\log(C/\pi) \approx 1.04$:
the cutoff terminates the sum after only the $n=1$ term,
but the prefactors $2\pi^2 n^4 e^{9u} - 3\pi n^2 e^{5u}$ can
contribute significantly even when the exponential suppression
is large.
At $u = 1.1$, the fixed-cutoff value of $\Phi$ differed from
the correct value by a relative error of $\sim 10^{-5}$.
This error propagated into the $L_4$ computation and produced
spurious negative values in the shift tail region.
The fix was to replace the cutoff rule with adaptive truncation
(termination when $|\text{term}/\text{sum}| < 10^{-(\dps+10)}$).
The regression tests at $u \in \{1.04, 1.1, 1.5\}$ verify that
the adaptive strategy is correct.

\emph{Per-script kernel copies.}
A second source of inconsistency was the practice of copying the
$\Phi$ implementation into each certificate script.
Minor differences between copies (rounding of constants, truncation
thresholds) led to situations where one script certified positivity
while another reported a failure at the same point.
The resolution was to centralize the kernel in a single module
(\texttt{certs/kernel.py}), imported by all scripts.
A runtime positivity guard (\texttt{Phi(u)\,$>$\,0} assertion)
provides an additional safety net.


\subsection{Lower-order total positivity}

\begin{lemma}[PF$_2$]\label{lem:PF2}
$K$ is PF$_2$.
\end{lemma}
\begin{proof}
By \Cref{lem:E0}, $g = -\log\Phi$ is strictly convex,
so $K(u) = e^{-g(|u|)}$ is strictly log-concave on $\R$,
hence PF$_2$ (see \cite[Ch.~4]{Karlin1968}).
\end{proof}

\begin{lemma}[PF$_3$]\label{lem:PF3}
$K$ is PF$_3$.
\end{lemma}
\begin{proof}
For every ordered triple $x_1 < x_2 < x_3$,
$y_1 < y_2 < y_3$, the Dodgson log-ratio at order~$3$ is
\[
  L_3 := \log(\Delta_{33}\,\Delta_{11}) - \log(\Delta_{31}\,\Delta_{13})
\]
where $\Delta_{ij}$ are the $2\times 2$ minors.
By PF$_2$, all $2\times 2$ minors are positive,
so $\det A_{3\times 3} > 0$ iff $L_3 > 0$.
Direct evaluation on a $5$-dimensional grid over
$[\delta, 1.0]^4 \times [-S, S]$
with $50{,}000$ evaluations at $60$-digit precision
yields $\min L_3 = 0.310$.
The shift tail follows from diagonal dominance
at large $|s|$;
coalescence from a Taylor expansion argument
analogous to the one developed for $L_4$ in \S\ref{sec:coalescence}.
\end{proof}


\subsection{Dodgson condensation}

The Desnanot--Jacobi identity~\cite{Dodgson1866,Bressoud1999}
expresses the $4\times 4$ determinant in terms of its
$3\times 3$ minors.
For $A = [K(x_i - y_j)]_{4\times 4}$,
let $\Delta_{ij}$ denote the $3\times 3$ minor obtained by deleting
row~$i$ and column~$j$, and let $\Delta_{ij,kl}$ denote the
$2\times 2$ minor obtained by deleting rows $i,k$ and columns $j,l$.

\begin{lemma}[Dodgson identity]\label{lem:dodgson}
\begin{equation}\label{eq:DJ}
  \det A \;\cdot\; \Delta_{14,41}
  \;=\; \Delta_{44}\,\Delta_{11} \;-\; \Delta_{41}\,\Delta_{14}.
\end{equation}
By PF$_3$, all $3\times 3$ minors $\Delta_{ij}$ are strictly positive.
By PF$_2$, the central $2\times 2$ minor $\Delta_{14,41}$ is
strictly positive.
Define the Dodgson log-ratio
\begin{equation}\label{eq:L4}
  L_4 := \log(\Delta_{44}\,\Delta_{11})
       - \log(\Delta_{41}\,\Delta_{14}).
\end{equation}
Then $\det A > 0$ if and only if $L_4 > 0$.
\end{lemma}
\begin{proof}
Rearranging~\eqref{eq:DJ}:
$\det A = (\Delta_{44}\Delta_{11} - \Delta_{41}\Delta_{14})/\Delta_{14,41}$.
Since $\Delta_{14,41} > 0$, the sign of $\det A$ equals the
sign of $\Delta_{44}\Delta_{11} - \Delta_{41}\Delta_{14}$,
which is positive iff $L_4 > 0$.
\end{proof}

The key advantage of working with $L_4$ rather than $\det A$
is \emph{numerical conditioning}.
The four $3\times 3$ minors are individually positive
(by PF$_3$), so $L_4$ is a difference of logarithms of
positive quantities.
In contrast, computing $\det A$ directly as a signed sum
of~$24$ terms incurs catastrophic cancellation:
at the Toeplitz minimum the determinant is of
order~$10^{-6}$ while the matrix entries are of order~$10^{-1}$.
The log-ratio $L_4$ is free from this cancellation and remains
well-conditioned throughout the parameter space.


% ═════════════════════════════════════════════════════════════════════════
\section{Geometry of the parameter space}\label{sec:geometry}
% ═════════════════════════════════════════════════════════════════════════

\subsection{Parametrization}

Translation invariance allows fixing $x_1 = 0$.
We parametrize the ordered $4$-tuples by gaps and a shift:
\begin{equation}\label{eq:param}
  x = (0,\; a_1,\; a_1{+}a_2,\; a_1{+}a_2{+}a_3), \qquad
  y = (s,\; s{+}b_1,\; s{+}b_1{+}b_2,\; s{+}b_1{+}b_2{+}b_3),
\end{equation}
with $a_k, b_k > 0$ and $s \in \R$.
The full parameter space is
$\Omega = (0,\infty)^6 \times \R$,
a seven-dimensional open domain.

\subsection{Partition and boundary parameters}

We fix three boundary parameters:
\begin{equation}\label{eq:params}
  \delta = 0.05, \qquad S = 0.4, \qquad G = 0.60.
\end{equation}
These define a partition of~$\Omega$ into four regions:

\begin{enumerate}[label=\textbf{R\arabic*}., leftmargin=2em]
  \item \textbf{Bounded interior.}
        All gaps $a_k, b_k \in [\delta, G]$ and $|s| \le S$.
        This is a compact $7$-dimensional box; its volume is the
        product of six gap ranges $[0.05, 0.60]$ and
        a shift range $[-0.4, 0.4]$.
        The grid minimum of $L_4$ is attained at the ``Toeplitz corner''
        $a_k = b_k = \delta$, $s = 0$; we certify a continuous
        neighborhood via a Taylor enclosure (\S\ref{sec:taylor}).
  \item \textbf{Large-gap tail.}
        At least one gap exceeds~$G$, with $|s| \le S$.
        The super-exponential decay of $\Phi$ makes the $4\times 4$
        matrix nearly block-diagonal, amenable to Schur complement
        analysis.
  \item \textbf{Large-shift tail.}
        $|s| > S$, arbitrary gaps $> 0$.
        For large $|s|$, the kernel values along the anti-diagonal
        of~$A$ are suppressed, and $L_4$ grows without bound.
  \item \textbf{Coalescence boundary.}
        At least one gap $a_k$ or $b_k$ lies in $(0, \delta)$.
        As a gap tends to zero, two rows (or columns) of~$A$ coalesce,
        and $\det A \to 0$ linearly with a positive leading coefficient.
\end{enumerate}

\begin{remark}[Overlap]\label{rem:overlap}
Region~\textbf{R1} and~\textbf{R4} overlap at the boundary
$a_k = \delta$ or $b_k = \delta$.
This overlap is by design: the dense grid certificate
(\Cref{lem:dense}) includes the value $a_k = \delta$,
providing a direct overlap check with the coalescence analysis.
Similarly, \textbf{R1} and~\textbf{R2} share the face $\max_k a_k = G$,
and the boundary grid (\Cref{lem:boundary}) covers gap values
up to and including~$G$.
\end{remark}

\subsection{Internal units}

All computations use two internal unit systems for exact integer
indexing:
\begin{itemize}[nosep]
  \item \emph{Centiunits} (cu): $1~\text{cu} = 0.01$ real.
        Used in the dense core and boundary grids,
        where the grid step is $1$~cu (dense) or $5$~cu (boundary).
  \item \emph{Half-centiunits} (hcu): $1~\text{hcu} = 0.005$ real.
        Used in the Taylor bridge, shift tail, and coalescence
        certificates, where finer resolution is required.
\end{itemize}
In internal units:
$\delta = 5$~cu $= 10$~hcu, $S = 40$~cu $= 80$~hcu,
$G = 60$~cu $= 120$~hcu.
All grid coordinates, loop bounds, and step sizes are integral
in the respective unit system.
Conversion to real numbers occurs only at JSON serialization;
the certification pipeline operates entirely in
\texttt{mpmath.mpf} arithmetic at $60$-digit precision.


% ═════════════════════════════════════════════════════════════════════════
\section{Interior positivity}\label{sec:interior}
% ═════════════════════════════════════════════════════════════════════════

\begin{lemma}[Dense core]\label{lem:dense}
On $[\delta, 0.13]^6 \times [-S, S]$ with grid step $0.01$:
$L_4 \ge 0.3077$. Zero failures in $2{,}657{,}205$ evaluations.
\end{lemma}
\begin{proof}
A $7$-dimensional grid with $9$ values per gap dimension
(from $0.05$ to $0.13$ in steps of $0.01$) and
$5$ shift values ($\{-0.4, -0.2, 0, 0.2, 0.4\}$)
produces $9^6 \times 5 = 2{,}657{,}205$ evaluation points.
At each point, $L_4$ is computed from four $3\times 3$ determinants
via the Dodgson log-ratio~\eqref{eq:L4}.
All evaluations return $L_4 > 0$, with a grid minimum of $0.3077$
achieved at the Toeplitz corner
$(a_k = b_k = \delta,\, s = 0)$.
\end{proof}

\begin{lemma}[Boundary extension]\label{lem:boundary}
On $[\delta, G]^6 \times [-S, S]$ with $\max(a_k, b_k) > 0.13$:
$L_4 \ge 0.647$. Zero failures in $587{,}925$ evaluations.
\end{lemma}
\begin{proof}
A non-uniform grid with $7$ gap values
($\{0.05, 0.10, 0.15, 0.20, 0.30, 0.45, 0.60\}$;
minimum step $0.05$, maximum step $0.15$)
and $5$ shift values covers the boundary region.
Configurations already covered by the dense core
($\max_k \le 0.13$) are skipped.
The remaining $587{,}925$ evaluations all satisfy $L_4 > 0$,
with a minimum of $0.647$.
\end{proof}

The grid certificates in Lemmas~\ref{lem:dense}
and~\ref{lem:boundary} establish $L_4 > 0$ at discrete points.
To extend the statement to the full continuous domain,
an interval-arithmetic Taylor model is used.
The next section develops this Taylor enclosure, which
is then applied adaptively across the bounded interior
(\Cref{cor:region1}).


% ═════════════════════════════════════════════════════════════════════════
\section{Taylor--Lipschitz certification near the minimum}\label{sec:taylor}
% ═════════════════════════════════════════════════════════════════════════

The Toeplitz corner $(a_k = b_k = \delta,\, s = 0)$ attains the
smallest value of~$L_4$ on the computational grid.
The grid value $L_4 = 0.3077$ is far from zero, but the grid
step ($0.01$ in the dense core) leaves open the possibility of
a negative excursion between grid points.
We develop the required interval-arithmetic second-order
Taylor enclosure on a small neighborhood of the minimum,
where certification is most demanding; the enclosure is then
applied adaptively across the full domain (\Cref{cor:region1}).

\subsection{The certification cell}

Define the \emph{mini-cell}
\[
  C = [\delta,\;\delta + h]^6 \times [-h_s,\; h_s]
\]
with $h = 0.025$ and $h_s = 0.05$ (in real units;
equivalently, gaps $\in [10, 15]$~hcu, shift $\in [-10, 10]$~hcu).
This cell encloses the Toeplitz minimum and overlaps with
the dense core grid, which includes the points
$a_k, b_k \in \{10, 12, 14\}$ hcu $= \{0.05, 0.06, 0.07\}$ real.

The cell~$C$ is subdivided along the shift coordinate into
$M_s = 8$ sub-boxes, each of the form
\[
  C_j = [\delta,\;\delta{+}h]^6 \times [s_j,\; s_j + \Delta s],
  \qquad \Delta s = 2h_s / M_s = 0.0125~\text{(real)}.
\]
The gap half-width is $r_g = h/2 = 0.0125$ and the shift
half-width is $r_s = \Delta s / 2 = 0.00625$.

\subsection{Taylor enclosure}

\begin{lemma}[Interval-arithmetic Taylor bridge]\label{lem:bridge}
$L_4(z) > 0$ for all $z \in C$.
The certified lower bound is $L_4 \ge 0.0724$.
\end{lemma}
\begin{proof}
All point evaluations of $L_4$ use \texttt{mpmath.iv} at $80$
digits with outward rounding, so every computed value is a
rigorous interval enclosing the true value.
For each sub-box~$C_j$ with center~$c$ and half-widths~$r$,
a second-order Taylor expansion gives
\begin{equation}\label{eq:taylor-bound}
  L_4(z) \;\ge\; L_4(c)
    - \sum_{k=1}^{7} G_k\, r_k
    - \tfrac{1}{2}\sum_{i,j=1}^{7} H_{ij}^{\mathrm{rig}}\, r_i\, r_j,
\end{equation}
where $G_k$ is the interval gradient envelope and
$H_{ij}^{\mathrm{rig}}$ is an interval-arithmetic upper bound
on the Hessian, both computed with outward rounding.
\end{proof}

The remainder of this section describes the construction of
$G_k$ and $H_{ij}^{\mathrm{rig}}$.

\subsection{Gradient envelope}

The true gradient $\nabla L_4(c)$ is approximated by central
differences.
A single step size provides no error control, so we employ
a \emph{self-consistency} approach with two step sizes
$\varepsilon_1 = 1$~hcu and $\varepsilon_2 = 2$~hcu.
Let $g_k^{(1)}$ and $g_k^{(2)}$ be the resulting central-difference
approximations for the $k$-th partial derivative.
Define
\begin{equation}\label{eq:grad-env}
  G_k \;:=\; \max\bigl(|g_k^{(1)}|,\,|g_k^{(2)}|\bigr)
          \;+\; |g_k^{(1)} - g_k^{(2)}|.
\end{equation}
The first term bounds the magnitude of the derivative estimate;
the second term bounds the truncation error.
Since the central-difference error is $O(\varepsilon^2)$,
the discrepancy
$|g_k^{(1)} - g_k^{(2)}| \approx |c_3|\cdot|\varepsilon_1^2 - \varepsilon_2^2|$
captures the leading error term, and $G_k$ provides an envelope
that covers the true $|\partial_k L_4(c)|$.
All point evaluations of~$g_k^{(j)}$ are performed in
\texttt{mpmath.iv}, so rounding errors are rigorously enclosed.
The residual $O(\varepsilon^4)$ truncation is at most
$\sim 10^{-12}$ for $\varepsilon \le 2$~hcu,
well below the $\ge 13\%$ headroom in the Hessian bound.

In practice, the gradient components are small near the Toeplitz
minimum (the minimum is a critical point in certain symmetry
directions), and the self-consistency discrepancies
$|g_k^{(1)} - g_k^{(2)}|$ are of order $10^{-4}$ --- four orders
of magnitude below the gradient envelope values.

\subsection{Rigorous Hessian bound}

The Hessian $H_{ij}(c) = \partial_i\partial_j L_4(c)$ is computed
at the sub-box center by second-order finite differences with
step $\varepsilon = 2$~hcu.
To bound $\sup_{z \in C_j} |H_{ij}(z)|$, we construct a
rigorous envelope:
\begin{equation}\label{eq:H-rig}
  H_{ij}^{\mathrm{rig}}
  \;=\; |H_{ij}(c)|
  \;+\; \sum_{k=1}^{7} |\partial_k H_{ij}(c)|\, r_k
  \;+\; 2\;\sum_{k=1}^{7} |\partial_k^2 H_{ij}(c)|\, r_k^2.
\end{equation}
The first-order correction $\partial_k H_{ij}$ (a third derivative
of~$L_4$) is estimated by central differences of the Hessian,
computed at the $14$ face-center points $c \pm r_k \hat{e}_k$.
The second-order correction $\partial_k^2 H_{ij}$ (a fourth
derivative) is estimated from the same face-center Hessian values
via the second-difference formula
$(H_+ - 2H_c + H_-)/r_k^2$.
The factor $2$ on the quadratic term provides rigorous headroom for
cross-partial contributions not captured by the diagonal
second-difference estimate; with the Taylor $\frac{1}{2}$ this gives
$4\times$ margin on the second-order remainder.

\subsection{Face-center validation}

The rigorous Hessian bound~\eqref{eq:H-rig} is validated at
all $15$ sample points (the center plus the $14$ face-centers).
For each sample~$p$ and each Hessian entry~$(i,j)$,
we verify $|H_{ij}(p)| \le H_{ij}^{\mathrm{rig}}$.
Across all $8$ sub-boxes:
\begin{itemize}[nosep]
  \item Zero violations (the ratio $|H_{ij}(p)| / H_{ij}^{\mathrm{rig}}$
        never exceeds~$1$).
  \item Maximum observed face-center ratio: $0.869$.
  \item The $\ge 13\%$ gap between the observed maximum and the bound
        provides margin against higher-order corrections.
\end{itemize}

\subsection{Certified lower bound}

Substituting into~\eqref{eq:taylor-bound}, the worst lower bound
across all $8$ sub-boxes is
\[
  \min_j \bigl[L_4(c_j) - \text{grad}_j - \text{hess}_j\bigr]
  \;=\; 0.0724,
\]
attained in the sub-box centered at $s = -1.25$~hcu
($s = -0.00625$ real).
The center value there is $L_4(c) = 0.538$,
the gradient correction is $0.275$, and the Hessian correction
is $0.190$.
Since $0.0724 > 0$, the Taylor enclosure certifies $L_4 > 0$
on the entire mini-cell~$C$.

\begin{remark}[Nature of the certificate]\label{rem:hessian-caveat}
The Taylor bridge is an \emph{interval-arithmetic certificate}:
every point evaluation of~$L_4$, its gradient, and its Hessian
uses \texttt{mpmath.iv} at $80$ digits with outward rounding,
following the approach used for the PF$_5$ counterexamples
(\S\ref{sec:pf5}).
The Hessian bound~$H_{ij}^{\mathrm{rig}}$ is verified at all
$14$ face-center sample points in interval arithmetic, with a
worst-case ratio of $0.87$ (i.e.\ $\ge 13\%$ headroom).
The Taylor model structure---center evaluation plus gradient and
Hessian bounds---is standard in computer-assisted proofs
(cf.~Tucker~\cite{Tucker2011}, Hales~\cite{Hales2005}).
\end{remark}

\begin{corollary}\label{cor:region1}
$L_4 > 0$ on $[\delta, G]^6 \times [-S, S]$.
\end{corollary}
\begin{proof}
The proof combines the adaptive tiling with the Schur
complement.  The Schur complement (\Cref{lem:schur}) certifies
$\det A > 0$ whenever \emph{any} gap exceeds $G_{\mathrm{Schur}} =
42$~hcu ($= 0.21$ real); see~\S\ref{sec:schur}.
It therefore suffices to prove $L_4 > 0$ on the sub-domain
where all six gaps are at most~$G_{\mathrm{Schur}}$.
We tile $[\delta,\,G_{\mathrm{tile}}]^6 \times [-S, S]$ with
$G_{\mathrm{tile}} = 44$~hcu ($2$~hcu overlap with the Schur threshold)
using an adaptive interval-arithmetic tiling in the style of
Tucker~\cite{Tucker2011} and Hales~\cite{Hales2005}.

For each box with center~$c$ and half-widths~$r$,
the second-order Taylor bound~\eqref{eq:taylor-bound} is
evaluated entirely in \texttt{mpmath.iv} at $80$~digits
with outward rounding.
The certification proceeds in two stages:
\begin{enumerate}[nosep]
  \item \emph{Simplified second-order}: center value
    $L_4(c)$, gradient envelope~$G_k$
    (\Cref{eq:grad-env}), and center Hessian $H(c)$
    with a safety factor of~$4$ on the quadratic term.
    This conservative bound covers
    the Hessian variation across the box
    (the face-center validation in \S\ref{sec:taylor}
    shows worst-case ratios $\le 0.87$ even with safety~$2$).
  \item \emph{Full Taylor model}: if the simplified bound
    fails, the full model with face-center Hessian
    verification (\Cref{lem:bridge}) is applied.
    If this also fails, the box is bisected along its
    longest dimension and both halves are re-certified.
\end{enumerate}
Symmetry reduces the computation: $L_4$ is invariant under
permutations of $(a_1, a_2, a_3)$ and of $(b_1, b_2, b_3)$
independently, and under $s \mapsto -s$.
This yields a factor of $72$ reduction.

The initial grid uses non-uniform spacing: step $2$~hcu
($= 0.01$ real) near the Toeplitz minimum
(gaps $\in [10, 22]$~hcu) where $L_4$ is smallest,
and step $\sim\!10$~hcu for gaps in $[22, 44]$~hcu.
The shift dimension uses step $10$~hcu.
All boxes certify at stage~1 (simplified second-order).
The worst certified lower bound is $L_4 \ge 5.0 \times 10^{-4}$,
achieved on a Toeplitz cell near $\text{gap} = 10$~hcu.

By the Dodgson identity (\Cref{lem:dodgson}),
$\det A > 0$ on the entire bounded interior.
\end{proof}


% ═════════════════════════════════════════════════════════════════════════
\section{Schur complement and large-gap regime}\label{sec:schur}
% ═════════════════════════════════════════════════════════════════════════

When at least one gap exceeds the threshold~$G = 0.60$,
the $4\times 4$ matrix $A = [K(x_i - y_j)]$ has a near
block-diagonal structure: the kernel $\Phi$ decays
super-exponentially, so entries coupling the ``far'' point
to the ``close cluster'' are exponentially suppressed.

\begin{lemma}[Schur complement]\label{lem:schur}
For any configuration with $|s| \le S$, all gaps $\ge \delta$,
and at least one gap $> G$: $\det A > 0$.
\end{lemma}
\begin{proof}
Let $g_k > G$ be the large gap.
Partition the four indices into a close cluster of three and
one far index.
Write
\[
  A = \begin{pmatrix} B & c \\ c^T & d \end{pmatrix},
\]
where $B$ is the $3\times 3$ submatrix of the close cluster,
$d = K(x_{\text{far}} - y_{\text{far}})$ is the diagonal entry of
the far point, and $c$ collects the cross terms.
By Sylvester's determinant formula,
$\det A = \det(B)\cdot(d - c^T B^{-1} c)$.

\emph{(i) $\det(B) > 0$}:
The close cluster has all gaps in $[\delta, G]$ with $|s| \le S$,
so $\det(B) > 0$ by PF$_3$ (\Cref{lem:PF3}), whose certificate
covers gaps up to~$1.0 \supset G = 0.60$.
For gaps exceeding $0.50$, the $3\times 3$ matrix is
diagonally dominant (off-diagonal-to-diagonal ratio $< 10^{-7}$)
by super-exponential decay of $\Phi$.

\emph{(ii) Schur ratio $\rho < 1$}:
Define $\rho := c^T B^{-1} c / d$.
Then $\det A > 0$ iff $\rho < 1$.
The ratio~$\rho$ is computed over all $6$ gap positions,
all $4$ possible Schur partitions (choosing which index is ``far''),
and $9$ shift values $s \in [-0.4, 0.4]$.
Results:

\begin{center}
\begin{tabular}{@{}cc@{}}
\toprule
$G$ & Worst $\rho$ \\
\midrule
0.30 & 0.990 \\
0.50 & 0.986 \\
0.60 & 0.985 \\
1.00 & 0.985 \\
\bottomrule
\end{tabular}
\end{center}

The ratio~$\rho$ is non-increasing in~$G$ for $G \ge 0.21$,
because the cross terms $c_j$ involve $\Phi$ evaluated at
distances $\ge G - 3\delta - S$, and $\Phi$ decays as
$\Phi(u) \sim \text{const} \cdot \exp(-\pi e^{4u})$.

\emph{Multiple large gaps.}
When two or more gaps exceed~$G$, the cross terms are further
suppressed.
The worst Schur ratio decreases with the number of large gaps:
$\rho \le 0.986$ for one or two large gaps,
$\rho \le 0.829$ for three or four,
and $\rho \le 0.024$ for five or more.
The certificate selects the best of the four possible Schur
partitions.

Therefore $\det A = \det(B)\cdot d\cdot(1 - \rho) > 0$.
\end{proof}

\begin{remark}\label{rem:schur-threshold}
Bisection on the Schur ratio over all positions and shift values
shows that $\rho < 1$ for $G \ge 0.21$.
We use $G = 0.60$ to ensure a comfortable margin and to keep
the bounded interior at a manageable size for the dense grid.
\end{remark}


% ═════════════════════════════════════════════════════════════════════════
\section{Shift tail regime}\label{sec:shift}
% ═════════════════════════════════════════════════════════════════════════

\begin{lemma}\label{lem:shift}
For $|s| > S = 0.4$ and all gaps $\ge \delta$:
$L_4 > 0$, hence $\det A > 0$.
\end{lemma}
\begin{proof}
By symmetry $L_4(\ldots, s) = L_4(\ldots, -s)$,
so it suffices to consider $s > 0$.

\emph{Minimizer.}
Let $f(s) = L_4(\delta,\ldots,\delta,s)$ denote $L_4$ at the
Toeplitz configuration (all gaps equal to~$\delta$).
A scan of $8{,}192$ gap configurations at $s = \pm S$ confirms
that the Toeplitz configuration achieves the minimum over all
gaps $\in \{\delta, 0.15, 0.30, G\}$: the minimum is $f(S) = 1.919$.

\emph{Bounded range $s \in [S, 2.0]$.}
The interval $[S, 2.0]$ is partitioned into $64$ sub-intervals
of width $\Delta s = 0.025$.
On each sub-interval $[s_k, s_{k+1}]$, a local Lipschitz bound gives
\[
  f(s) \;\ge\; f_{\min} - \ell_k \cdot \Delta s / 2,
\]
where $f_{\min} = \min(f(s_k), f(s_{k+1}))$ and the local
Lipschitz constant is
\[
  \ell_k = \max\bigl(|f'(s_k)|,\,|f'(s_{k+1})|\bigr)
         + \max\bigl(|f''(s_k)|,\,|f''(s_{k+1})|\bigr)\cdot\Delta s/2.
\]
Here $f'$ and $f''$ are computed by central differences
at the sub-interval endpoints.
Worst certified margin: $1.8096$ (at $s = S$, where $f = 1.9190$
and the Lipschitz correction is $0.1095$).
All $64$ intervals pass with positive margin.

The following table summarizes the key certified values:
\begin{center}
\begin{tabular}{@{}lrl@{}}
\toprule
Quantity & Value & Source \\
\midrule
$f(S) = L_4(\delta,\ldots,\delta, S)$ & $1.9190$ & grid eval \\
Worst margin over all intervals & $1.8096$ & Lipschitz cert \\
Gap configurations tested at $s = \pm S$ & $8{,}192$ & gap scan \\
$f(2.0)$ & $1127.5$ & grid eval \\
\bottomrule
\end{tabular}
\end{center}

\emph{Asymptotic range $s > 2.0$.}
On the entire grid from $s = S$ to $s = 2.0$,
the derivative $f'(s)$ is strictly positive:
$f$ increases monotonically from $1.919$ to $1127.5$.
Analytically, $L_4 \sim 4\,g''(s)\,\delta^2$ as $s \to \infty$,
where $g'' > 0$ by \Cref{lem:E0} and $g''(s) \to 16\pi^2 e^{8s}$.
Therefore $L_4 \to \infty$.

Since $f'(s) > 0$ for all $s \ge S$ on the grid and $f \to \infty$,
the function~$f$ is increasing beyond the grid,
completing the certificate.
\end{proof}


% ═════════════════════════════════════════════════════════════════════════
\section{Coalescence boundary}\label{sec:coalescence}
% ═════════════════════════════════════════════════════════════════════════

\begin{lemma}\label{lem:coal}
For any configuration with at least one gap
$a_k \in (0, \delta]$ (other gaps $\ge \delta$, $|s| \le S$):
$\det A > 0$.
\end{lemma}
\begin{proof}
As $a_k \to 0$, rows~$i$ and~$i{+}1$ of~$A$ coalesce
and $\det A \to 0$.
The determinant vanishes linearly:
\begin{equation}\label{eq:coal-expansion}
  \det A \;=\; a_k \cdot \Ntilde \;+\; O(a_k^2),
\end{equation}
where $\Ntilde := \lim_{a_k \to 0} \det A / a_k$ is the
first-order Taylor coefficient.
Geometrically, $\Ntilde$ is a mixed minor involving $K'$,
positive by the PF$_3$ structure of~$K$.

\emph{Leading coefficient.}
$\Ntilde$ is computed by Richardson extrapolation at
$h = 1$~hcu and $h = 0.5$~hcu:
$\Ntilde_{\text{extrap}} = 2\Ntilde(h/2) - \Ntilde(h)$,
which eliminates the $O(h)$ term and gives a second-order
accurate estimate.
Across all $6$ gap positions and
$s \in \{0, \pm 0.2, \pm 0.4\}$:
$\min \Ntilde = 8.248 \times 10^{-23} > 0$.

\emph{Direct evaluation.}
The determinant $\det A$ is evaluated directly for single
small gaps ranging from $0.1$ to $10$~hcu
($0.0005$ to $0.05$ real)
across all $6$ positions and $9$ shift values
($432$ evaluations, $\min \det A = 1.494\times 10^{-23}$).
For two simultaneously small gaps:
$1{,}200$ evaluations across all $\binom{6}{2} = 15$ pairs,
all positive ($\min \det A = 1.418\times 10^{-26}$).

The certified values are summarized below:
\begin{center}
\begin{tabular}{@{}lrl@{}}
\toprule
Quantity & Value & Method \\
\midrule
$\min \Ntilde$ (leading coeff.) & $8.248\times 10^{-23}$ & Richardson \\
$\min \det A$ (single small gap) & $1.494\times 10^{-23}$ & direct eval ($432$ pts) \\
$\min \det A$ (two small gaps) & $1.418\times 10^{-26}$ & direct eval ($1{,}200$ pts) \\
$\min D_4/\text{gap}$ & $1.494\times 10^{-22}$ & ratio scan \\
\bottomrule
\end{tabular}
\end{center}

\emph{Overlap with Region~1.}
At gap $= \delta$
(the boundary with the bounded interior):
$\det A(\delta,\ldots,\delta,0) = 3.81\times 10^{-6} > 0$,
consistent with the grid value $L_4 = 0.3077$ at the same point.
\end{proof}

\begin{remark}
The small magnitudes ($10^{-23}$ to $10^{-26}$) arise because
$\det A$ is a $4\times 4$ determinant of kernel values that are
themselves products of exponentials.
At $60$-digit precision, these values are resolved with
over $30$ significant digits of margin.
\end{remark}


% ═════════════════════════════════════════════════════════════════════════
\section{Assembly of the proof}\label{sec:assembly}
% ═════════════════════════════════════════════════════════════════════════

\begin{proof}[Proof of \Cref{thm:main}]
We show $\det[K(x_i - y_j)]_{4\times 4} > 0$ for all
ordered $x_1 < x_2 < x_3 < x_4$ and
$y_1 < y_2 < y_3 < y_4$.
By translation invariance, fix $x_1 = 0$
and use the gap-shift parametrization~\eqref{eq:param}.

\begin{enumerate}[label=(\roman*)]
\item \emph{Bounded interior}
  (gaps $\in [\delta, G]$, $|s| \le S$):
  $L_4 > 0$ by \Cref{cor:region1} (adaptive interval-arithmetic
  tiling; worst certified lower bound $L_4 \ge 5.0 \times 10^{-4}$).
  By the Dodgson identity (\Cref{lem:dodgson}), $\det A > 0$.

\item \emph{Large-gap tail}
  ($\exists$ gap $> G$, $|s| \le S$):
  Schur complement (\Cref{lem:schur});
  $\rho < 0.986$, hence $\det A > 0$.

\item \emph{Large-shift tail} ($|s| > S$):
  Local Lipschitz bounds and monotonicity (\Cref{lem:shift});
  worst margin $1.8096$, hence $\det A > 0$.

\item \emph{Coalescence} ($\exists$ gap $< \delta$):
  Positive leading Taylor coefficient
  $\Ntilde \ge 8.25 \times 10^{-23}$ (\Cref{lem:coal}),
  with overlap verification at gap~$= \delta$.
\end{enumerate}
The four regions cover $\Omega = (0,\infty)^6 \times \R$.
\end{proof}


% ═════════════════════════════════════════════════════════════════════════
\section{Failure of PF\texorpdfstring{$_5$}{5}}\label{sec:pf5}
% ═════════════════════════════════════════════════════════════════════════

\begin{proof}[Proof of \Cref{prop:not-pf5}]
For the Toeplitz configuration with parameters
$(u_0, h) = (0.01, 0.05)$,
interval-arithmetic computation at $80$-digit precision with
outward rounding (\texttt{mpmath.iv}) gives
\[
  \det[K(u_0 + (i-j)h)]_{5\times 5}
  \;\in\; [-1.8473,\, -1.8472] \times 10^{-9},
\]
a rigorous enclosure entirely below zero.
In total, $12$ distinct counterexample configurations have been
verified, with $u_0 \in \{0.01, 0.02, 0.03\}$ and
$h \in [0.03, 0.05]$
(\texttt{certs/verify\_pf5.py}, output frozen in
\texttt{results/pf5.json}).
\end{proof}

Together with \Cref{thm:main}, this shows that PF$_4$ is
the sharp total-positivity order of the de~Bruijn--Newman kernel.
The failure at order~$5$ arises in a specific parameter regime
(small, equally spaced gaps near~$0$),
suggesting that the PF$_4$ property is not far from a borderline
phenomenon.


% ═════════════════════════════════════════════════════════════════════════
\section{Computational framework and reproducibility}\label{sec:comp}
% ═════════════════════════════════════════════════════════════════════════

\subsection{Architecture}

The companion repository\footnote{%
  \url{https://github.com/ScypyonX/pf4-dodgson-global-proof}.}
is organized around a single canonical kernel implementation
(\texttt{certs/kernel.py}).
All certificates import from this module; no per-script copies
of~$\Phi$ or its derivatives exist.
Shared parameters ($\delta$, $S$, $G$, precision, series truncation
limit) are centralized in \texttt{certs/config.py}.
Each certificate is produced by a single deterministic command
and outputs a JSON file to \texttt{results/}, recording the
result, precision, and environment metadata (Python version,
platform, \texttt{mpmath} version, and all domain parameters).

\subsection{Certificates at a glance}

Table~\ref{tab:certs} maps each proof component to its certificate
script, domain, and frozen output.

\begin{table}[ht]
\caption{Certificate summary: each row corresponds to one
proof component, its computational domain, the script that
produces it, and the frozen JSON output.}\label{tab:certs}
\centering
\small
\begin{tabular}{@{}llllll@{}}
\toprule
Region & Script & Domain & Evals & Result & JSON \\
\midrule
K2+E0 & \texttt{master\_cert} & $[0,1]$, 1D & 501 &
  $\Phi'<0$, $g''\!>\!0$ & \texttt{master.json} \\
PF$_3$ & \texttt{master\_cert} & $[\delta,1.0]^4\!\times\![-S,S]$ & 50\,k &
  $L_3 \ge 0.310$ & \texttt{master.json} \\
\textbf{R1} core & \texttt{master\_cert} &
  $[\delta,0.13]^6\!\times\![-S,S]$ & 2.66\,M &
  $L_4 \ge 0.3077$ & \texttt{master.json} \\
\textbf{R1} bdy & \texttt{master\_cert} &
  $[\delta,G]^6\!\times\![-S,S]$ & 588\,k &
  $L_4 \ge 0.647$ & \texttt{master.json} \\
\textbf{R1} bridge & \texttt{boundary\_taylor} &
  $[\delta,\delta\!+\!h]^6\!\times\![-h_s,h_s]$ & $8\!\times\!15$ &
  $L_4 \ge 0.0724$ (ia) & \texttt{boundary\_taylor.json} \\
\textbf{R1} tiling & \texttt{continuous\_tiling} &
  $[\delta,G_{\mathrm{tile}}]^6\!\times\![-S,S]$ & adaptive &
  $L_4 \ge 5.0 \times 10^{-4}$ (ia) & \texttt{continuous\_tiling.json} \\
\textbf{R2} & \texttt{schur\_tail} &
  $\exists$gap$\!>\!G$ & bisect &
  $\rho < 0.986$ & \texttt{schur\_tail.json} \\
\textbf{R3} & \texttt{shift\_tail} &
  $|s|\!\in\![S,2.0]$ & 64+8\,k &
  margin $\ge 1.810$ & \texttt{shift\_tail.json} \\
\textbf{R4} & \texttt{coalescence} &
  $\exists$gap$\!<\!\delta$ & 1\,632 &
  $\Ntilde \ge 8.25\!\times\!10^{-23}$ & \texttt{coalescence.json} \\
PF$_5$ & \texttt{verify\_pf5} &
  Toeplitz & 12 &
  $\det_5 < 0$ (ia) & \texttt{pf5.json} \\
\bottomrule
\end{tabular}
\end{table}

\subsection{Computational summary}

\begin{center}
\begin{tabular}{@{}llrlr@{}}
\toprule
Certificate & Dim & Evaluations & Precision & Result \\
\midrule
K2 + E0              & 1D  & 501           & 60 dig & $\Phi'<0$, $g''>0$ \\
PF$_3$               & 5D  & 50,000        & 60 dig & $L_3 \ge 0.310$ \\
Dense core           & 7D  & 2,657,205     & 60 dig & $L_4 \ge 0.3077$ \\
Boundary             & 7D  & 587,925       & 60 dig & $L_4 \ge 0.647$ \\
Taylor bridge        & 7D  & $8\times 15$  & 80 dig (ia) & $L_4 \ge 0.0724$ \\
Adaptive tiling      & 7D  & adaptive      & 80 dig (ia) & $L_4 \ge 5.0 \times 10^{-4}$ \\
Schur ratio          & 7D  & bisection     & 60 dig & $\rho < 0.986$ \\
Shift tail           & 1D  & 64 + 8,192    & 60 dig & margin $\ge 1.810$ \\
Coalescence          & 7D  & 1,632         & 60 dig & $\Ntilde \ge 8.25\!\times\!10^{-23}$ \\
PF$_5$ falsification & 5D  & 12 configs    & 80 dig (ia) & $\det_5 < 0$ \\
\bottomrule
\end{tabular}
\end{center}

Total: over $3.3 \times 10^6$ kernel evaluations for \Cref{thm:main}
(including the PF$_3$ prerequisite),
with grid certificates at $60$-digit multiprecision
(\texttt{mpmath.mp.dps\,=\,60}).
The Taylor bridge, adaptive tiling, and PF$_5$ falsification
use interval arithmetic (\texttt{mpmath.iv})
at $80$ digits with outward rounding.

\subsection{Reproducibility}

\begin{center}
\begin{tabular}{@{}lll@{}}
\toprule
Command & Certificate & Time \\
\midrule
\texttt{python certs/pf4\_master\_certificate.py}  & K2+E0, PF$_3$, R1, Schur    & $\sim$3\,min \\
\texttt{python certs/pf4\_boundary\_taylor.py}     & Taylor bridge (mini-cell)    & $\sim$30\,s \\
\texttt{python certs/pf4\_continuous\_tiling.py}   & Adaptive tiling (R1)         & $\sim$hours \\
\texttt{python certs/pf4\_schur\_tail.py}          & Schur ratio                  & $\sim$10\,s \\
\texttt{python certs/pf4\_shift\_tail.py}          & Shift tail                   & $\sim$20\,s \\
\texttt{python certs/pf4\_coalescence.py}          & Coalescence                  & $\sim$15\,s \\
\texttt{python certs/verify\_pf5.py}               & PF$_5$ falsification         & $\sim$1\,s \\
\texttt{python tests/test\_kernel.py}              & Unit tests (12 tests)        & $\sim$5\,s \\
\midrule
\texttt{bash verify\_all.sh}                      & All of the above             & $\sim$1\,h \\
\bottomrule
\end{tabular}
\end{center}

Frozen JSON outputs in \texttt{results/} serve as reference
reports: each file records the numerical result together with
a metadata block containing the Python version, platform,
\texttt{mpmath} version, precision, series truncation limit,
and all domain parameters.
JSON values are serialized at IEEE\,754 double precision;
certification decisions (sign tests, inequality checks)
are performed in \texttt{mpmath.mpf} at $60$ digits for grid
certificates, and in \texttt{mpmath.iv} at $80$ digits for the
Taylor bridge, adaptive tiling, and PF$_5$ falsification,
before serialization.
Any discrepancy between a fresh run and the frozen output can
be traced to environment differences via this metadata.
Requirements: Python~$\ge 3.10$, \texttt{mpmath}~$\ge 1.3$.

\subsection{Limitations}

All components of the proof either use exact evaluations at grid
points combined with analytical arguments (Schur monotonicity,
diagonal dominance, coalescence Taylor expansion), or are
interval-certified: the PF$_5$ counterexamples (\S\ref{sec:pf5}),
the Taylor bridge (\Cref{lem:bridge}), and the adaptive tiling
(\Cref{cor:region1}) all use
\texttt{mpmath.iv} with outward rounding at $80$ digits.
The interval-arithmetic components use a Taylor model
rather than direct interval evaluation of~$L_4$ on boxes,
because the Dodgson determinant computation exhibits
catastrophic cancellation under naive interval extension
in~$7$~dimensions.


% ═════════════════════════════════════════════════════════════════════════
\section{Discussion and outlook}\label{sec:discussion}
% ═════════════════════════════════════════════════════════════════════════

The result PF$_4$ (and the failure of PF$_5$) gives a precise
characterization of the total-positivity order of the
de~Bruijn--Newman kernel.
We conclude with several directions for future work.

\emph{Direct interval evaluation.}
The Taylor bridge uses an interval-arithmetic Taylor model
because direct interval evaluation of the Dodgson log-ratio
on~$7$-dimensional boxes suffers from exponential dependency
blowup in the determinant computation.
A reformulation using preconditioned determinants or
alternative total-positivity criteria might enable direct
interval evaluation, eliminating the Taylor model layer.

\emph{Dependence on $\Lambda$.}
The PF$_4$ property is established for the kernel $K$ at
$\lambda = 0$ (the physical case).
For $\lambda > 0$, the kernel $K_\lambda$ is a Gaussian
perturbation of~$K$, and total positivity is expected to
increase (Gaussian kernels are TP$_\infty$).
Conversely, for $\lambda < 0$ the total-positivity order may decrease.
A systematic study of PF$_r(K_\lambda)$ as a function of $\lambda$
could shed light on the structure of the deformation.

\emph{Higher-order analysis.}
The PF$_5$ counterexample occurs at small, equally spaced gaps
near the origin.
Understanding why the $5\times 5$ determinant changes sign in
this regime --- and whether the sign change is related to
the structure of the theta-series coefficients --- may yield
analytic insight into the borderline between PF$_4$ and the
failure of PF$_5$.

\emph{Implications for zero distributions.}
The PF$_4$ property implies a variation-diminishing bound
of order~$3$ for convolutions with~$K$.
The precise implications for the zero distribution of
$H_\lambda(z)$ at $\lambda = 0$ merit further investigation,
particularly in connection with the de~Bruijn--Newman constant.


% ═════════════════════════════════════════════════════════════════════════
% References
% ═════════════════════════════════════════════════════════════════════════

\begin{thebibliography}{99}

\bibitem{Bressoud1999}
D.\,M.~Bressoud,
\emph{Proofs and Confirmations: The Story of the Alternating Sign
Matrix Conjecture},
Cambridge Univ.\ Press, 1999.

\bibitem{CNV1986}
G.~Csordas, T.\,S.~Norfolk, R.\,S.~Varga,
\emph{The Riemann hypothesis and the Tur\'an inequalities},
Trans.\ Amer.\ Math.\ Soc.\ \textbf{296} (1986), 521--541.

\bibitem{deBruijn1950}
N.\,G.~de~Bruijn,
\emph{The roots of trigonometric integrals},
Duke Math.\ J.\ \textbf{17} (1950), 197--226.

\bibitem{Dodgson1866}
C.\,L.~Dodgson,
\emph{Condensation of determinants},
Proc.\ Roy.\ Soc.\ London \textbf{15} (1866), 150--155.

\bibitem{FallatJohnson2011}
S.\,M.~Fallat, C.\,R.~Johnson,
\emph{Totally Nonnegative Matrices},
Princeton Univ.\ Press, 2011.

\bibitem{Hales2005}
T.\,C.~Hales,
\emph{A proof of the Kepler conjecture},
Ann.\ of Math.\ (2) \textbf{162} (2005), 1065--1185.

\bibitem{Karlin1968}
S.~Karlin,
\emph{Total Positivity}, Vol.~I,
Stanford Univ.\ Press, 1968.

\bibitem{Newman1976}
C.\,M.~Newman,
\emph{Fourier transforms with only real zeros},
Proc.\ Amer.\ Math.\ Soc.\ \textbf{61} (1976), 245--251.

\bibitem{Pinkus2010}
A.~Pinkus,
\emph{Totally Positive Matrices},
Cambridge Univ.\ Press, 2010.

\bibitem{PlattTrudgian2021}
D.\,J.~Platt, T.\,S.~Trudgian,
\emph{The de Bruijn--Newman constant is at most $0.2$},
J.\ Reine Angew.\ Math.\ (to appear), arXiv:2104.12285.

\bibitem{RodgersTao2020}
B.~Rodgers, T.~Tao,
\emph{The de~Bruijn--Newman constant is non-negative},
Forum Math.\ Pi \textbf{8} (2020), e6.

\bibitem{Schoenberg1951}
I.\,J.~Schoenberg,
\emph{On P\'olya frequency functions. I},
J.\ Analyse Math.\ \textbf{1} (1951), 331--374.

\bibitem{Tucker2011}
W.~Tucker,
\emph{Validated Numerics: A Short Introduction to Rigorous Computations},
Princeton Univ.\ Press, 2011.

\end{thebibliography}

\end{document}
